\section*{\zihao{-2}致谢}
\addcontentsline{toc}{section}{致谢}

\textit{I\~{} feel something so right}

\textit{Doing the wrong thing}

\textit{I\~{} feel something so wrong}

\textit{Doing the right thing}

\textit{I couldn't lie, couldn't lie, couldn't lie}

\textit{Everything that kills me makes me feel alive}

\hfill \textit{------ ``Counting Stars'' OneRepublic}

\medskip

首先，我要衷心感谢我的两位导师——中国科学技术大学的陈雪老师与美国约翰霍普金斯大学的李昕老师。本论文的相当一部分工作基于我大三期间在两位导师指导下合作完成的英文论文 \textit{Explicit Min-wise Hash Families with Optimal Size} \cite{xue_shengtang_xin_min_wise}，该论文已被理论计算机领域顶级会议 the 37th Annual ACM-SIAM Symposium on Discrete Algorithms (SODA 2026) 接收。同时，也感谢陈雪老师在本中文毕业论文写作过程中给予的耐心指导与细致修改。

陈雪老师是我进入理论计算机科学领域的引路人。自大一暑假参加老师组织的伪随机讨论班以来，我逐渐对这一研究随机性与结构性之间关系的方向产生了浓厚兴趣。在此后的学习与研究过程中，陈雪老师不仅在学术上悉心指导我阅读经典文献、探索研究问题，也在个人发展与升学规划方面给予了极大的支持与帮助，使我受益匪浅。同时，我也感谢陈雪老师在我担任大二春季学期其《算法基础》课程助教期间的指导与支持。

李昕老师是我在大三暑假与大四秋季学期期间于约翰霍普金斯大学访问交流时的导师。在此期间，李昕老师在多个研究课题上给予了深入指导，使我的科研能力得到进一步提升。同时，李昕老师也在未来发展与学术道路选择方面给予了宝贵建议。

另外，我要感谢中国科学技术大学的陈士祥老师。我在大三秋季学期担任其《运筹学》课程助教期间，陈士祥老师在助教工作上给予了耐心指导；在我后续升学过程中，陈士祥老师亦为我撰写了推荐信，提供了宝贵支持。

此外，我还要感谢美国加州大学戴维斯分校的 Slobodan Mitrovi\'{c} 老师。在我大二暑假期间初次前往海外交流时，Slobodan 老师在科研与生活两方面都给予了我诸多帮助与关照，并包容了我蹩脚的英语，使我顺利度过了那段重要的适应阶段。

我也感谢所有曾给予我指导与帮助的前辈和同学们。感谢 Slobodan 老师的博士生許文宏学长，在加州期间在生活与学业上给予我的帮助；感谢李昕老师的博士生钟妍学姐与毛松涛学长，在巴尔的摩期间的指导，并在学术规划方面提供了许多宝贵建议；感谢中科大的周翟恩和同学，作为少有的理论计算机方向同伴，你敏捷而深刻的思考常常给予我启发，衷心祝愿你未来在美国的求学生活顺利且光明。

接下来，我要把最温柔最绵长的感谢，留给一路同行的朋友与始终在我身后的家人。

我要感谢那些贯穿我学生时代的朋友们——从初高中到大学，几近十年的同窗岁月，与你们真挚的友情是我求学生涯中最真实、也最令人眷恋的部分。感谢中科大的洪博远同学、魏海亮同学、张子悦同学、黄骏坤同学，以及几位后辈李存淳学妹、陈睿涵学弟、陈天添学弟；感谢北京大学的俞畅同学、清华大学的林而立同学与谢濡键同学。一起并肩学习、讨论问题，一起吃饭、闲聊、出游，一起搞抽象，一起哈哈大笑——与你们一起度过的时光，总是明亮而鲜活。

我要感谢我的父母，黄强汉先生与柯荔娜女士。感谢你们始终以最坚定最无私的方式支持着我。是你们给予我自由选择道路的勇气，让我可以在并不容易的学术道路上走得从容而笃定；是你们在我面临压力迷茫和疲惫的时候，给予我最朴素却最有力量的安慰，和我并肩面对与解决问题；是你们在我每次回到家中时都准备着温暖的床铺与可口的饭菜，告诉我“家是你最温暖的港湾”；是你们带我来到这大千世界，让我健康成长，展翅翱翔。感谢你们的理解、包容与爱，这是我一切尝试与坚持的底气所在。

最后，我想把最特别的一句感谢，留给你——我最最可爱的猫猫 capoo 小姐。与你相遇，是我一生的幸运，感谢你的陪伴。

% 最后，感谢我自己，感谢自己虽然会迷茫会焦虑，却一直慢慢的往前走。

\centering \textit{We'll be counting stars!}

\medskip

\hfill{2026 年 5 月于合肥}